% Bagean: metodhe
% Bab: bukti
% Pérangan: miwiti bukti

\documentclass[../../../include/open-logic-section]{subfiles}

\begin{document}

\olfileid{mth}{prf}{str}

\olsection{Miwiti Bukti}

Nanging saka ngendi kudu miwiti?

Panjenengan diwènèhi pratelan kang kudu dibuktèkaké, mula pratelan
iku prayoga dadi bab pungkasan kang kasebut ing bukti. Mesthi,
ing wiwitan panjenengan bisa \emph{ngumumaké} yèn arep mbuktekaké
pratelan iku, nanging aja nganggep pratelan iku minangka anggepan.
Tulisen perkara kang arep dibuktèkaké ing pérangan ngisor kertas
anyar, supaya tujuané tansah kèlingan.

Sabanjuré, mbokmenawa ana sawatara anggepan kang bisa digunakaké.
Bab iki bakal luwih cetha nalika ing pérangan sabanjuré kita ngrembug
\emph{jinis} bukti. Tulisen anggepan-anggepan iku ing pérangan
ndhuwur kertas lan tandhanana kanthi cetha yèn iku anggepan:
upamané, yèn nganggep $p$, tulisen ``anggepen $p$'' utawa
``supaya $p$''. Pungkasané, mbokmenawa ana dhéfinisi ing pitakonan
kang kudu dingertèni. Panjenengan bisa didhawuhi migunakaké
dhéfinisi tartamtu, utawa bisa uga ana sawatara dhéfinisi ing
anggepan utawa kesimpulan kang arep digayuh.
\emph{Tulisen dhéfinisi-dhéfinisi iku lan mesthèkna yèn panjenengan
ngerti tegesé.}

Cara nata bukti uga gumantung marang wujud pitakonan. Pérangan
sabanjuré njlentrehaké carané nata bukti miturut jinis ukarané.

\end{document}
